% !TEX TS-program = xelatex
% !TEX encoding = UTF-8 Unicode
\documentclass[11pt]{article}
\usepackage[margin=1in]{geometry}

\usepackage{xcolor}
\usepackage{hyperref}
\hypersetup{
  colorlinks=true,
  citecolor=violet
}

\usepackage[citestyle=authoryear, bibstyle=apa]{biblatex}
\addbibresource{references.bib}

\usepackage{fontspec} 
\setmainfont{Charis SIL}


\title{Proposal to L\&C 45th Annual Gender Studies Symposium}
\author{Wren Arbor (she/her), Student of Linguistics at Reed College (wren.arbor@reed.edu)}
\date{November 7th, 2025}

\begin{document}
\maketitle

\noindent \textbf{Title of Presentation:} Talking while Trans: Transfeminine Navigations of the Cisgender Listener

\noindent \textbf{Format of Presentation:} Individual Paper


\section{Brief Biographical Information}

My name is Wren Arbor (she/her), and I am currently an undergraduate majoring in linguistics at Reed College. My area of interest is trans linguistics, especially trans sociophonetics, a subfield of linguistics dedicated to understanding the relationship between gender and the sounds of the voice from a trans perspective. I am a white American middle-class queer trans  woman in her early twenties, and my work stems from this positionality. The L\&C Gender Studies symposium would be my first time giving a conference talk, and I'm glad to have heard it is a supportive and welcoming venue for this! I am lucky to have studied under professors who have conducted research on language and gender, including my thesis advisor Allison Casar (see their \href{https://sites.google.com/view/allisontrcasar/research-projects?authuser=0#h.tnc9nc20jtg2}{manuscripts in preparation}), and my professors Kara Becker and Sameer Khan (see \cite{becker_beyond_2022}). 

\section{Proposal}

As a senior undergraduate at Reed College, I am working on a thesis which will be complete by the end of the 2025-26 academic year. I propose here to present this research at the L\&C 45th Annual Gender Studies Symposium.

My thesis aims to contribute to our understanding of identity,  gender performance, the transgender experience, and transphobia by investigating how transfeminine individuals change their pronunciation of the /s/ sound (as in the first sound in `sand' or the last sound in `face') based on the trans status of their conversation partner. 

Scholars of language and gender have come to understand the sounds of the voice as resources for gender performativity (\cite{bucholtz_feminist_2014}). Increasingly, research has looked to trans voices and perspectives for insight on this process, investigating the role of the voice in presentations of gender identities beyond the binary (see \cite{zimman_transgender_2020}). This research typically involves the interpretation of connections between speakers' identities and the sounds of their voices. These connections are only interpretable within a set of culturally-available associations (see \cite{calder_whose_2021}) that a listener might apply in perception of a speaker's voice. Existing studies often choose to focus on only one set of associations, and thus only one type of theoretical listener-- at present, either the set of associations of the speakers' own local queer community (eg. \cite{calder_sissy_2019}), or the hegemonic set of normative associations as represented by findings among the cis straight white population (eg. \cite{hazenberg_language_2012, podesva_sexuality_2015}). The latter type of study, by interpreting trans voices in their deviance from or compliance with cisnormativity, runs the risk of tacitly enforcing these standards, marginalizing anti-normative trans voice practices (\cite{calder_interrogating_2025}). Furthermore, focus within a single type of theoretical listener -- hegemonic or otherwise -- often sidelines the intersubjective nature of identity, by which identity emerges in the establishment and contrast of self and other, speaker and listener (see \cite{bucholtz_locating_2010}). Recent research has taken steps to place more even focus on both sides of the speaker-listener pair, investigating how trans voices are percieved by listeners (\cite{jessee_cisgender_2025}) and trans people's understandings of these perceptions (\cite{heffernan_constructing_2023}). However, research which treats the listener as a variable, investigating how trans speakers' understandings of how they will be percieved manifest in their vocal practice, has not yet been conducted. By filling this gap, I hope to demonstrate the diverse strategies trans speakers employ to resist transphobic perceptions.

Language and gender research has found that the pronunciation of the /s/ sound is subject to Western hegemonic white cisheteronormative gender norms; higher-pitched /s/ pronunciations are expected of and associated with feminity (see \cite{calder_gay_2020}). \textcite{jessee_cisgender_2025} find that cisgender listeners impose this standard more strictly on trans voices than on cis voices, expecting /s/ pitch to be higher from trans women than cis women and lower for trans men than cis men. Although transfeminine people are likely not consciously aware of /s/ pitch as an index of femininity, they nonetheless seem to understand that cis listeners apply strict gender norms to their voices, and have strategies to navigate these expectations. One of the transfeminine speakers in \textcite{heffernan_constructing_2023} noted that when ``interacting with the cis universe'' she finds herself ``deliberately modulating [her] voice'' to make it ``more girly'', and that this is ``a kind of survival tactic... to try and avoid any possibility that [she]'ll get misgendered'' (23). From this context, it seems possible that an increase in /s/ pitch is part of the vocal modulation she is speaking of. I have chosen to investigate /s/ pitch for this reason. 

In order to investigate how trans speakers' understandings of how their voices are percieved manifest in their vocal practice, I will audio-record transfeminine speakers having one conversation with a cis stranger, and one conversation with a trans stranger. I will then compare each speakers' /s/ pitch between the two interactions. One possible outcome would be in line with the experience of Heffernan's speaker, where speakers raise their /s/ pitch to accomodate the demands of cisnormativity and avoid the reprocussions of choosing not to do so. However, it is likely that some transfeminine people feel differently than the speaker in Heffernan, instead choosing not to raise /s/ pitch as a move of resistance or to present a uniquely trans feminine identity beyond cis norms. For these reasons, I hypothesize that trans speakers will employ diverse vocal strategies which may help them defend their self-determined identities, avoid misgendering, resist gender norms, or create new forms of gender presentation.

\textbf{Anticipated AV needs:}
My presentation would be a google slide deck, so I would need to be able to put that up on a screen.

\printbibliography
\end{document}